\documentclass[12pt]{article}

\usepackage[a4paper, margin=2.5cm]{geometry}
\usepackage{amsmath}
\usepackage{amssymb}
\usepackage{amsthm}
\usepackage[colorlinks=true, linkcolor=blue, citecolor=blue, urlcolor=blue]{hyperref}
\usepackage{booktabs}
\usepackage{natbib}

\newtheorem{theorem}{Theorem}
\newtheorem{proposition}[theorem]{Proposition}
\newtheorem{corollary}[theorem]{Corollary}
\newtheorem{definition}[theorem]{Definition}
\newtheorem{remark}[theorem]{Remark}
\newtheorem{example}[theorem]{Example}

\newcommand{\Fcal}{\mathcal{F}}
\newcommand{\VW}{V_{\mathrm{W}}}
\newcommand{\HW}{H_{\mathrm{W}}}
\newcommand{\Sadm}{\mathcal{S}_{\mathrm{adm}}}
\newcommand{\Rst}{R_{\mathrm{st}}}
\newcommand{\dUrel}{\delta U_{\mathrm{rel}}}

\title{The Closure Residual:\\Nonlinear Remainder\\Beyond the Witten-Hamiltonian Regime}
\author{%
  Lee Smart\\[4pt]
  \textit{Institute of Vibrational Field Dynamics}\\[2pt]
  \texttt{contact@vibrationalfielddynamics.org}\\[2pt]
  \texttt{@vfd\_org}%
}
\date{April 2026}

\begin{document}

\maketitle

\noindent\textbf{Status:} Preprint (not peer-reviewed)

\begin{abstract}
Paper~XVIII derived an exact nonlinear closure-Nelson wave equation and showed that the nonlinear residual vanishes exactly at the Paper~XIV equilibrium, so that the wave equation reduces to the Schr\"odinger equation with Witten Hamiltonian $\HW$; for perturbations around equilibrium, this linear Schr\"odinger dynamics persists to leading order. The remaining discrepancy away from equilibrium is a state-dependent nonlinear term with coefficient $\sigma^2$ in the Schr\"odinger-normalised form.

The present paper studies this \textit{closure residual} directly. Rewriting the dynamics in equilibrium-centred normal form $i\partial_t\Psi = \HW\Psi + \dUrel[\Psi]\Psi$, the residual is shown to vanish exactly at the Paper~XIV equilibrium and to measure curvature deviation of the probability amplitude from equilibrium structure. Structural properties are derived: global phase invariance, modulus-only dependence, locality, and near-equilibrium smallness. The full nonlinear equation is shown to conserve the $L^2$ norm exactly, confirming that probability is preserved even in the nonlinear regime. A gauge analysis establishes that the residual cannot be removed by phase redefinition or linear unitary transformation, and that amplitude-based absorption changes the physical content of the theory.

A dimensionless residual-strength parameter $\Xi$ is defined, providing a regime classifier for the dynamics: a linear Schr\"odinger regime ($\Xi \ll 1$), a mixed regime ($\Xi \sim 1$), and a nonlinear closure-dominated regime ($\Xi \gg 1$). A quadratic worked example demonstrates each regime explicitly.

The paper does not interpret the residual as a modification of quantum mechanics. It establishes it as a mathematically well-defined object --- the closure-sector remainder after linear quantum behaviour has been recovered --- whose physical interpretation is deferred to subsequent work.
\end{abstract}

%==================================================================
\section{Introduction}\label{sec:intro}
%==================================================================

Paper~XVIII~\citep{SmartXVIII} established that forward/backward pairing of closure stochastic processes produces a nonlinear Schr\"odinger-type wave equation:
\begin{equation}\label{eq:XVIII}
    i\sigma^2\partial_t\Psi = -\frac{\sigma^4}{2}\Delta\Psi + \left(-\sigma^2\VW + \sigma^4\frac{\Delta|\Psi|}{|\Psi|}\right)\Psi
\end{equation}
where $\VW = |\nabla\Fcal|^2/(2\sigma^2) - \Delta\Fcal/2$ is the Witten potential. At the Paper~XIV~\citep{SmartXIV} equilibrium and to leading order near equilibrium, this reduces to the linear Schr\"odinger equation $i\partial_t\Psi = \HW\Psi$ with Witten Hamiltonian $\HW = -(\sigma^2/2)\Delta + \VW$.

The sole remaining obstruction to global linear Schr\"odinger dynamics is the state-dependent term $\sigma^4\Delta|\Psi|/|\Psi|$ (or equivalently $\sigma^2\Delta|\Psi|/|\Psi|$ after normalisation to Schr\"odinger form). Paper~XVIII identified this as a structurally derived residual rather than a defect of the construction. The present paper takes the next step: it isolates, classifies, and analyses this residual mathematically.

\subsection*{Scope}

This paper does not interpret the residual as new physics. It establishes the residual's mathematical properties so that subsequent work can determine its physical status. The paper answers: \textit{what kind of object is the residual, and when does it matter?}

%==================================================================
\section{Equilibrium-Centred Normal Form}\label{sec:normalform}
%==================================================================

Define the stationary amplitude $\Rst = \sqrt{\rho_{\mathrm{st}}} \propto e^{-\Fcal/\sigma^2}$, where $\rho_{\mathrm{st}} = Z^{-1}e^{-2\Fcal/\sigma^2}$ is the Paper~XIV equilibrium density. The key identity (derived in Paper~XVIII) is:
\begin{equation}\label{eq:equil_identity}
    \frac{\Delta\Rst}{\Rst} = \frac{|\nabla\Fcal|^2}{\sigma^4} - \frac{\Delta\Fcal}{\sigma^2} = \frac{2\VW}{\sigma^2}
\end{equation}

Dividing~\eqref{eq:XVIII} by $\sigma^2$ and using~\eqref{eq:equil_identity}, the closure-Nelson wave equation takes the \textit{equilibrium-centred normal form}:

\begin{equation}\label{eq:normalform}
    \boxed{i\,\partial_t\Psi = \HW\Psi + \dUrel[\Psi]\,\Psi}
\end{equation}

where the \textit{closure residual} is:
\begin{equation}\label{eq:residual}
    \boxed{\dUrel[\Psi] = \sigma^2\left(\frac{\Delta|\Psi|}{|\Psi|} - \frac{\Delta\Rst}{\Rst}\right)}
\end{equation}

\begin{remark}
This form separates the dynamics into two components: the linear Witten-Hamiltonian evolution $\HW\Psi$ (which is exact Schr\"odinger), and the nonlinear closure residual $\dUrel[\Psi]\Psi$ (which vanishes at equilibrium). The residual measures how far the current amplitude profile deviates from the equilibrium curvature structure. The coefficient $\sigma^2$ results from dividing Paper~XVIII's wave equation (which has coefficient $\sigma^4$ in the undivided form $i\sigma^2\partial_t\Psi = \ldots$) by $\sigma^2$ to obtain the Schr\"odinger-normalised form.
\end{remark}

%==================================================================
\section{Structural Properties of the Residual}\label{sec:structure}
%==================================================================

\begin{proposition}[Equilibrium vanishing]\label{prop:vanishing}
$\dUrel[\Psi_{\mathrm{st}}] = 0$, where $\Psi_{\mathrm{st}} = \Rst$ is the equilibrium amplitude.
\end{proposition}

\begin{proof}
Direct substitution: $|\Psi_{\mathrm{st}}| = \Rst$, so $\Delta|\Psi_{\mathrm{st}}|/|\Psi_{\mathrm{st}}| = \Delta\Rst/\Rst$, and the difference vanishes.
\end{proof}

\begin{proposition}[Phase invariance]\label{prop:phase}
$\dUrel[e^{i\theta}\Psi] = \dUrel[\Psi]$ for any constant $\theta \in \mathbb{R}$. More generally, $\dUrel[\Psi]$ depends only on $|\Psi|$ and its spatial derivatives, not on the phase of $\Psi$.
\end{proposition}

\begin{proof}
$|e^{i\theta}\Psi| = |\Psi|$, so $\Delta|e^{i\theta}\Psi|/|e^{i\theta}\Psi| = \Delta|\Psi|/|\Psi|$.
\end{proof}

\begin{proposition}[Locality]\label{prop:locality}
On any open set where $|\Psi| > 0$, $\dUrel[\Psi](\psi)$ depends only on $|\Psi|$, $\nabla|\Psi|$, and $\Delta|\Psi|$ evaluated at $\psi$. It is a local functional of the modulus.
\end{proposition}

\begin{proof}
On $\{|\Psi| > 0\}$: $\Delta|\Psi|/|\Psi| = \Delta\log|\Psi| + |\nabla\log|\Psi||^2$, which involves only local derivatives of $|\Psi|$.
\end{proof}

\begin{proposition}[Near-equilibrium scaling]\label{prop:scaling}
For $\Psi = \Psi_{\mathrm{st}}(1 + \varepsilon\eta)$ with $\varepsilon \ll 1$ and $\eta$ smooth:
\begin{equation}
    \dUrel[\Psi] = O(\varepsilon)
\end{equation}
More precisely, to leading order:
\begin{equation}\label{eq:linearised_residual}
    \dUrel[\Psi] = \varepsilon\sigma^2\left(\Delta\mathrm{Re}(\eta) + 2\nabla\log\Rst\cdot\nabla\mathrm{Re}(\eta)\right) + O(\varepsilon^2)
\end{equation}
\end{proposition}

\begin{proof}
Write $R = |\Psi| = \Rst(1 + \varepsilon\,\mathrm{Re}(\eta)) + O(\varepsilon^2)$. Then:
\begin{align}
    \frac{\Delta R}{R} &= \frac{\Delta\Rst}{\Rst} + \varepsilon\left(\Delta\mathrm{Re}(\eta) + 2\nabla\log\Rst\cdot\nabla\mathrm{Re}(\eta)\right) + O(\varepsilon^2)
\end{align}
Subtracting $\Delta\Rst/\Rst$ gives~\eqref{eq:linearised_residual}.
\end{proof}

\begin{proposition}[Characterisation of vanishing]\label{prop:vanishing_class}
$\dUrel[\Psi] = 0$ if and only if $|\Psi|$ satisfies:
\begin{equation}
    \frac{\Delta|\Psi|}{|\Psi|} = \frac{\Delta\Rst}{\Rst} = \frac{2\VW}{\sigma^2}
\end{equation}
That is, the residual vanishes precisely for states whose amplitude has the same curvature-to-amplitude ratio as the equilibrium profile.
\end{proposition}

\begin{proof}
Direct from definition~\eqref{eq:residual}.
\end{proof}

%==================================================================
\section{Conservation Laws}\label{sec:conservation}
%==================================================================

\begin{proposition}[Norm conservation]\label{prop:norm}
The full nonlinear equation~\eqref{eq:normalform} preserves the $L^2$ norm:
\begin{equation}
    \frac{d}{dt}\int|\Psi|^2\,d\psi = 0
\end{equation}
\end{proposition}

\begin{proof}
Write~\eqref{eq:normalform} as $i\partial_t\Psi = \HW\Psi + \dUrel\Psi$. Then:
\begin{align}
    \frac{d}{dt}\int|\Psi|^2 &= 2\,\mathrm{Re}\int\Psi^*\,\partial_t\Psi\,d\psi \notag \\
    &= 2\,\mathrm{Re}\int\Psi^*\left(-i\HW\Psi - i\,\dUrel\Psi\right)d\psi
\end{align}
For the Witten term: $\mathrm{Re}\int\Psi^*(-i\HW\Psi) = -\mathrm{Im}\langle\Psi, \HW\Psi\rangle = 0$ since $\HW$ is self-adjoint and $\langle\Psi,\HW\Psi\rangle$ is real.

For the residual term: $\mathrm{Re}\int\Psi^*(-i\,\dUrel\Psi) = -\mathrm{Im}\int\dUrel\,|\Psi|^2 = 0$ since $\dUrel[\Psi]$ is \textit{real-valued} (it depends only on modulus derivatives via~\eqref{eq:residual}).
\end{proof}

\begin{remark}[Physical significance]
Norm conservation holds \textit{exactly} for the full nonlinear equation, not just in the linearised regime. Probability is preserved even when the residual is large. This is because $\dUrel$ acts as a real multiplicative potential: it rotates the phase of $\Psi$ but does not drain or amplify its modulus. The closure-Nelson dynamics is therefore probability-conserving in all regimes.
\end{remark}

\begin{remark}[On unitarity vs.\ norm conservation]
The closure dynamics preserves the $L^2$ norm but is not generated by a skew-adjoint operator. In this sense it is norm-preserving but not unitary in the standard quantum-mechanical sense: the generator $\HW + \dUrel[\Psi]$ is self-adjoint only in the linearised regime where $\dUrel \to 0$.
\end{remark}

\subsection*{Continuity equation}

Since the norm is conserved, the density $\rho = |\Psi|^2$ satisfies a continuity equation. Writing $\Psi = R\,e^{iS/\sigma^2}$, the Madelung decomposition gives:
\begin{equation}
    \partial_t\rho + \nabla\cdot(\rho\,\nabla S) = 0
\end{equation}
which is exact for the full nonlinear equation (inherited from Paper~XVIII's Madelung pair).

\subsection*{Energy functional}

Define the candidate energy:
\begin{equation}\label{eq:energy}
    E[\Psi] = \langle\Psi, \HW\Psi\rangle = \int\left(\frac{\sigma^2}{2}|\nabla\Psi|^2 + \VW|\Psi|^2\right)d\psi
\end{equation}

\begin{proposition}[Energy non-conservation]\label{prop:energy}
The Witten energy $E[\Psi]$ is \textbf{not} conserved by the full nonlinear dynamics. Its time derivative is:
\begin{equation}
    \frac{dE}{dt} = -2\,\mathrm{Im}\int\dUrel[\Psi]\,\Psi^*\,\HW\Psi\,d\psi
\end{equation}
which vanishes at equilibrium and to leading order near equilibrium, but is generically nonzero.
\end{proposition}

\begin{proof}
$dE/dt = 2\,\mathrm{Re}\langle\partial_t\Psi, \HW\Psi\rangle = 2\,\mathrm{Re}\langle -i\HW\Psi - i\dUrel\Psi, \HW\Psi\rangle$. The first term gives $-2\,\mathrm{Im}\langle\HW\Psi, \HW\Psi\rangle = 0$ (since $\HW$ is self-adjoint). The second gives $-2\,\mathrm{Im}\langle\dUrel\Psi, \HW\Psi\rangle = -2\,\mathrm{Im}\int\dUrel[\Psi]\,\Psi^*\,\HW\Psi\,d\psi$, which is generically nonzero since $\dUrel$ is real but $\Psi^*\HW\Psi$ is generally complex.
\end{proof}

\begin{remark}
The functional $E[\Psi] = \langle\Psi, \HW\Psi\rangle$ is the \textit{linear} Witten energy --- the energy of the linearised Schr\"odinger regime. It is not necessarily the conserved energy of the full nonlinear theory. Norm is conserved but linear Witten energy is not: the nonlinearity exchanges energy between the Witten-Hamiltonian sector and the residual sector. Whether a modified energy functional incorporating the residual exists and is conserved remains open.
\end{remark}

%==================================================================
\section{Gauge and Removability Analysis}\label{sec:gauge}
%==================================================================

We test whether the residual can be absorbed by admissible transformations.

\subsection{Phase redefinition}

\begin{proposition}[Non-removability by phase]\label{prop:phase_gauge}
No transformation $\Psi \to \Psi\,e^{iG(|\Psi|)}$ with $G$ depending only on $|\Psi|$ can remove $\dUrel$.
\end{proposition}

\begin{proof}
Under $\Psi \to \Psi\,e^{iG}$, the modulus $|\Psi|$ is unchanged, so $\dUrel[\Psi\,e^{iG}] = \dUrel[\Psi]$. The residual is phase-invariant (Proposition~\ref{prop:phase}).
\end{proof}

\subsection{Linear unitary transformation}

\begin{proposition}[Non-removability by linear unitary]\label{prop:unitary_gauge}
No time-independent unitary transformation $\Psi \to U\Psi$ (with $U$ a fixed linear operator) can remove $\dUrel$ from the equation.
\end{proposition}

\begin{proof}
Under $\Psi \to U\Psi$, the equation transforms as $i\partial_t(U\Psi) = U\HW U^\dagger(U\Psi) + \dUrel[U\Psi](U\Psi)$. The residual $\dUrel[U\Psi] = \sigma^2(\Delta|U\Psi|/|U\Psi| - \Delta\Rst/\Rst)$ depends on $|U\Psi|$, which generically differs from $|\Psi|$ for non-trivial $U$. The residual is modified but not eliminated.
\end{proof}

\subsection{Equilibrium-amplitude factorisation}

The substitution $\Psi = \Rst\,\Phi$ (used in Paper~XVI to eliminate the drift) transforms the equation. Since $|\Psi| = \Rst|\Phi|$:
\begin{equation}
    \frac{\Delta|\Psi|}{|\Psi|} = \frac{\Delta(\Rst|\Phi|)}{\Rst|\Phi|}
\end{equation}

Expanding: $\Delta(\Rst|\Phi|) = |\Phi|\Delta\Rst + 2\nabla\Rst\cdot\nabla|\Phi| + \Rst\Delta|\Phi|$, so:
\begin{equation}
    \frac{\Delta|\Psi|}{|\Psi|} = \frac{\Delta\Rst}{\Rst} + \frac{2\nabla\log\Rst\cdot\nabla|\Phi|}{|\Phi|} + \frac{\Delta|\Phi|}{|\Phi|}
\end{equation}

The residual becomes:
\begin{equation}\label{eq:residual_phi}
    \dUrel = \sigma^2\left(\frac{2\nabla\log\Rst\cdot\nabla|\Phi|}{|\Phi|} + \frac{\Delta|\Phi|}{|\Phi|}\right)
\end{equation}

\begin{remark}
The equilibrium piece $\Delta\Rst/\Rst$ cancels exactly, but the residual is not eliminated: it is re-expressed in terms of $|\Phi|$ and its derivatives. The residual vanishes if and only if $|\Phi|$ is constant (i.e.\ $|\Psi| \propto \Rst$). The factorisation simplifies the residual's dependence but does not remove it.
\end{remark}

\subsection{Summary of removability}

\begin{table}[ht]
\centering
\caption{Removability of the closure residual under candidate transformations.}
\label{tab:removability}
\begin{tabular}{@{}lll@{}}
\toprule
Transformation & Effect on $\dUrel$ & Removable? \\
\midrule
Global phase $e^{i\theta}$ & Invariant & No \\
Phase redefinition $e^{iG(|\Psi|)}$ & Invariant & No \\
Linear unitary $U$ & Modified, not eliminated & No \\
Equilibrium factorisation $\Rst\Phi$ & Re-expressed in $|\Phi|$ & No \\
Nonlinear amplitude transform & May absorb, changes theory & Open \\
\bottomrule
\end{tabular}
\end{table}

The residual is robust under all transformations that preserve the physical content of the theory. It can only be absorbed by transformations that alter the equation's structure.

\begin{proposition}[Homogeneity under amplitude rescaling]\label{prop:homogeneity}
For any constant $\lambda > 0$:
\begin{equation}
    \dUrel[\lambda\Psi] = \dUrel[\Psi]
\end{equation}
Thus the residual depends on the shape of the modulus profile, not on its overall normalisation.
\end{proposition}

\begin{proof}
Since $|\lambda\Psi| = \lambda|\Psi|$, $\Delta|\lambda\Psi|/|\lambda\Psi| = \Delta|\Psi|/|\Psi|$. The equilibrium term is unchanged, so $\dUrel[\lambda\Psi] = \dUrel[\Psi]$.
\end{proof}

%==================================================================
\section{Residual Strength and Dynamical Regimes}\label{sec:regimes}
%==================================================================

Define the dimensionless \textit{residual-strength parameter}:
\begin{equation}\label{eq:Xi}
    \Xi[\Psi](\psi) = \frac{|\dUrel[\Psi](\psi)|}{|\VW(\psi)| + \sigma^2|\Delta\Psi/\Psi|/2}
\end{equation}
as a pointwise ratio of the residual magnitude to the characteristic scale of the Witten-Hamiltonian terms.

Three regimes:

\begin{table}[ht]
\centering
\caption{Dynamical regimes classified by residual strength.}
\label{tab:regimes}
\begin{tabular}{@{}lll@{}}
\toprule
Regime & Condition & Dynamics \\
\midrule
Linear Schr\"odinger & $\Xi \ll 1$ & $i\partial_t\Psi \approx \HW\Psi$ \\
Mixed & $\Xi \sim 1$ & Witten + residual compete \\
Nonlinear closure & $\Xi \gg 1$ & Residual dominates \\
\bottomrule
\end{tabular}
\end{table}

\begin{remark}
The linear Schr\"odinger regime contains the equilibrium ($\Xi = 0$) and its perturbative neighbourhood. The nonlinear closure regime corresponds to configurations whose density profiles differ strongly from $\rho_{\mathrm{st}}$ --- for instance, highly localised or rapidly oscillating amplitude profiles. The mixed regime is the transitional zone where both components contribute comparably. This pointwise quantity is defined on regions where $|\Psi| > 0$; global (integrated) versions and treatment near nodal sets are left for future work.
\end{remark}

%==================================================================
\section{Worked Example: Quadratic Closure Functional}\label{sec:example}
%==================================================================

\begin{example}[Quadratic closure]\label{ex:quadratic}
Let $\Fcal(\psi) = \frac{1}{2}k\psi^2$ in one dimension.

\subsection*{Equilibrium structure}

$\Rst = \mathcal{N}\,e^{-k\psi^2/(2\sigma^2)}$, $\VW = k^2\psi^2/(2\sigma^2) - k/2$, and $\Delta\Rst/\Rst = 2\VW/\sigma^2 = k^2\psi^2/\sigma^4 - k/\sigma^2$.

The Witten Hamiltonian has spectrum $E_n = nk$ with eigenstates $\phi_n$.

\subsection*{Residual for Witten eigenstates}

For $\Psi = \phi_n$ (an eigenstate of $\HW$), $|\Psi| = |\phi_n|$, and:
\begin{equation}
    \dUrel[\phi_n] = \sigma^2\left(\frac{\Delta|\phi_n|}{|\phi_n|} - \frac{k^2\psi^2}{\sigma^4} + \frac{k}{\sigma^2}\right)
\end{equation}

For the ground state $\phi_0 \propto e^{-k\psi^2/(2\sigma^2)} = \Rst$: $\dUrel[\phi_0] = 0$. $\checkmark$

For the first excited state $\phi_1 \propto \psi\,e^{-k\psi^2/(2\sigma^2)}$: $|\phi_1| = |\psi|\,\mathcal{N}\,e^{-k\psi^2/(2\sigma^2)}$, and one computes (away from $\psi = 0$):
\begin{equation}
    \frac{\Delta|\phi_1|}{|\phi_1|} = \frac{k^2\psi^2}{\sigma^4} - \frac{3k}{\sigma^2} + \frac{2}{\psi^2}
\end{equation}

So: $\dUrel[\phi_1] = \sigma^2(2/\psi^2 - 2k/\sigma^2) = 2\sigma^2/\psi^2 - 2k$.

The residual diverges as $\psi \to 0$ (where $|\phi_1|$ has a node) and is negative for large $|\psi|$. The nodal divergence is expected whenever $\Delta|\Psi|/|\Psi|$ is evaluated at zeros of $|\Psi|$; the expression is understood on regions where $|\Psi| > 0$, and nodal sets require separate treatment. The residual is nonzero for excited eigenstates --- confirming that only the ground state lies exactly in the Schr\"odinger regime.

\subsection*{Residual for a shifted Gaussian}

Consider a Gaussian shifted from equilibrium: $|\Psi| = \mathcal{N}\,e^{-k(\psi - a)^2/(2\sigma^2)}$ for some displacement $a > 0$. Then:
\begin{equation}
    \frac{\Delta|\Psi|}{|\Psi|} = \frac{k^2(\psi-a)^2}{\sigma^4} - \frac{k}{\sigma^2}
\end{equation}

The residual is:
\begin{equation}
    \dUrel = \sigma^2\left(\frac{k^2(\psi-a)^2}{\sigma^4} - \frac{k}{\sigma^2} - \frac{k^2\psi^2}{\sigma^4} + \frac{k}{\sigma^2}\right) = \frac{k^2}{\sigma^2}\left[(\psi-a)^2 - \psi^2\right] = \frac{k^2(-2a\psi + a^2)}{\sigma^2}
\end{equation}

This is a linear function of $\psi$ (plus a constant) --- nonzero, but finite and smooth. The residual strength grows linearly with displacement $a$, confirming that near-equilibrium perturbations produce small residuals.

\subsection*{Regime summary}

\begin{table}[ht]
\centering
\caption{Residual behaviour for different states in the quadratic model.}
\label{tab:states}
\begin{tabular}{@{}lll@{}}
\toprule
State & $\dUrel$ & Regime \\
\midrule
Ground state $\phi_0$ & $0$ & Exact Schr\"odinger \\
Shifted Gaussian (small $a$) & $O(a)$ & Linear Schr\"odinger \\
First excited state $\phi_1$ & Diverges at nodes & Nonlinear closure \\
Highly localised state & Large & Nonlinear closure \\
\bottomrule
\end{tabular}
\end{table}

\end{example}

%==================================================================
\section{Interpretation}\label{sec:interpretation}
%==================================================================

The closure residual $\dUrel[\Psi]$ has been shown to be:
\begin{itemize}
    \item zero at equilibrium (Proposition~\ref{prop:vanishing}),
    \item phase-invariant (Proposition~\ref{prop:phase}),
    \item local and modulus-dependent (Proposition~\ref{prop:locality}),
    \item small near equilibrium (Proposition~\ref{prop:scaling}),
    \item not removable by phase, unitary, or equilibrium-factorisation transformations (Section~\ref{sec:gauge}),
    \item probability-conserving (Proposition~\ref{prop:norm}).
\end{itemize}

These properties establish the residual as a mathematically well-defined, robust, non-trivial remainder relative to linear Schr\"odinger dynamics. It encodes the curvature mismatch between the actual amplitude profile and the equilibrium amplitude profile.

\begin{remark}[On the role of the residual]
The residual should not be interpreted as a failure to recover Schr\"odinger dynamics. Rather, it is the exact remainder that distinguishes the full nonlinear closure dynamics from its equilibrium tangent limit. Its presence is therefore expected: it encodes the deviation of the system from the equilibrium regime in which linear quantum behaviour is recovered.
\end{remark}

\subsection*{What the residual is not}

\begin{itemize}
    \item It is not a numerical artefact: it is an exact consequence of the closure-Nelson pairing.
    \item It is not a gauge-removable quantity: it survives all standard transformations.
    \item It is not a defect of the construction: it is the closure-sector remainder after linear quantum behaviour has been recovered.
\end{itemize}

\subsection*{What the residual may be}

Three interpretive possibilities remain open:
\begin{enumerate}
    \item \textbf{A physical nonlinear correction.} The residual may represent a genuine departure from linear quantum mechanics in regimes far from closure equilibrium.
    \item \textbf{An effective-theory artefact.} The residual may be an artefact of the closure approximation that would vanish in a more complete framework.
    \item \textbf{A structurally meaningful quantity.} The residual may encode physical information about the closure landscape that is not captured by the linearised Schr\"odinger dynamics.
\end{enumerate}

Determining which of these holds requires further analysis --- in particular, whether the residual produces observable consequences or can be constrained by known experimental bounds. This is deferred to subsequent work.

%==================================================================
\section{Discussion and Limitations}\label{sec:discussion}
%==================================================================

\subsection{Exact results}

\begin{itemize}
    \item The equilibrium-centred normal form~\eqref{eq:normalform} is exact.
    \item The structural properties (Propositions~\ref{prop:vanishing}--\ref{prop:vanishing_class}) are proved.
    \item Norm conservation (Proposition~\ref{prop:norm}) holds for the full nonlinear equation.
    \item The gauge non-removability results (Propositions~\ref{prop:phase_gauge}--\ref{prop:unitary_gauge}) are proved.
    \item The residual for the quadratic example is computed explicitly.
\end{itemize}

\subsection{Regime-specific results}

\begin{itemize}
    \item The near-equilibrium scaling $\dUrel = O(\varepsilon)$ (Proposition~\ref{prop:scaling}) quantifies smallness.
    \item The residual-strength parameter $\Xi$ classifies regimes but is defined pointwise; a global (integral) version would be more robust.
\end{itemize}

\subsection{Open questions}

\begin{itemize}
    \item Does a modified energy functional (incorporating the residual) exist that is conserved?
    \item Can the residual be bounded in terms of a distance from equilibrium (e.g.\ relative entropy)?
    \item Does the nonlinear equation have special solutions not present in the linear theory?
    \item Can the residual be constrained by existing experimental bounds on nonlinear quantum mechanics?
\end{itemize}

%==================================================================
\section{Conclusion}\label{sec:conclusion}
%==================================================================

Paper~XIX isolates and classifies the closure residual --- the sole remaining discrepancy within the closure-Nelson formulation between the paired-process dynamics and the linear Schr\"odinger equation.

\textbf{What is established.} The residual $\dUrel[\Psi] = \sigma^2(\Delta|\Psi|/|\Psi| - \Delta\Rst/\Rst)$ is:
\begin{itemize}
    \item exactly zero at equilibrium,
    \item perturbatively small near equilibrium,
    \item phase-invariant, local, and modulus-dependent,
    \item not removable by standard gauge transformations,
    \item probability-conserving in all regimes.
\end{itemize}

\textbf{What this means for the programme.} Paper~XVIII solved the kinetic problem: the imaginary kinetic term emerges from the forward/backward pairing. Paper~XIX classifies the remaining nonlinear residual and shows it is a robust, non-trivial structural quantity. The dynamics is approximately linear Schr\"odinger in the equilibrium-centred regime ($\Xi \ll 1$) and genuinely nonlinear when the residual is not negligible.

\textbf{Programme arc.}
\begin{itemize}
    \item Papers~XII--XIV: dissipative closure dynamics,
    \item Paper~XV: phase and interference,
    \item Paper~XVI: evolution equation and structural gap,
    \item Paper~XVII: dissipative obstruction theorem,
    \item Paper~XVIII: Nelson pairing and near-equilibrium Schr\"odinger recovery,
    \item Paper~XIX: classification of the closure residual.
\end{itemize}

The next stage is to determine the physical status of the residual: whether it is an observable nonlinear correction, an effective-theory artefact, or a structural signature of the closure landscape. That determination will decide whether the closure framework extends, modifies, or exactly recovers standard quantum mechanics.

%==================================================================
\bibliographystyle{plainnat}
\bibliography{references}

\end{document}
